\chapter{Math Background}
\label{app1}


\section{Fr{\'e}chet Differentiability and Derivative}\label{section-frechet-differentiable}

We provide details on the Fr{\'e}chet differentiability and derivative. Throughout this section, we let $\calH$ be a real Hilbert space, $J: \calH \to \Real$ be a map, $\calB \parens{\calH, \Real}$ denote the collection of all bounded linear operators from $\calH$ to $\Real$, and, similarly, $\calB \parens{\calH, \calH}$ denote the collection of all bounded linear operators from $\calH$ to itself. All materials of this section come from Section 2.6 in \textcite{Bauschke2017-he} and Section 5.1 in \textcite{Denkowski2013-ke}. 

\begin{definition}[Fr{\'e}chet differentiability and derivative]\label{def-frechet-differentiable}
	The map $J$ is said to be \textit{(first-order) Fr{\'e}chet differentiable} at $f \in \calH$ if there exists an operator $\Diff J \parens{f} \in \calB \parens{\calH, \Real}$ such that 
	\begin{align}\label{eq-frederivative}
	    \lim_{\substack{\norm{g}_{\calH} \to 0 \\ g \neq 0}} \frac{\abs[\big]{J \parens{f + g} - J \parens{f} - \Diff J \parens{f}\parens{g}}}{\norm{g}_{\calH}} = 0, 
	\end{align}
	and the operator $\Diff J \parens{f}$ is called the \textit{(first-order) Fr{\'e}chet derivative}. The map $J$ is said to be \textit{(first-order) Fr{\'e}chet differentiable on $\calH$} if it is Fr{\'e}chet differentiable at all $f \in \calH$. 
\end{definition}

\begin{proposition}
	Suppose $J$ is Fr{\'e}chet differentiable at $f \in \calH$ and the Fr{\'e}chet derivative $\Diff J \parens{f}$ exists. Then, $\Diff J \parens{f}$ is unique. 
\end{proposition}

\begin{remark}
	If $J$ is Fr{\'e}chet differentiable at $f \in \calH$, we then can write 
	\begin{align}\label{eq-frechet-derivative-interpretation}
		J \parens{f + g} = J \parens{f} + \Diff J \parens{f} \parens{g} + o \parens{\norm{g}_{\calH}}, 
	\end{align}
	for all $g \in \calH$ in a small neighborhood of the origin, where $o \parens{\norm{g}_{\calH}}$ denotes $\frac{o \parens{\norm{g}_{\calH}}}{\norm{g}_{\calH}} \to 0$ as $\norm{g}_{\calH} \to 0$. Thus, from \eqref{eq-frechet-derivative-interpretation}, we see $J \parens{f} + \Diff J \parens{f} \parens{g}$ provides the best linear approximation of $J$ in a small neighborhood of $f$, which is the similar interpretation of the derivative of a real-valued function of a single variable. 
\end{remark}

Frech{\'e}t derivative shares many properties of the derivative of a real-valued function of a single variable. The following proposition lists two properties we use in studying the Frech{\'e}t differentiability and deriving the Frech{\'e}t derivative of the log-partition functional $A$ in Chapter \ref{ch2}. 

\begin{proposition}\label{prop-property-frechet-differentiable}
\leavevmode
	\begin{enumerate}[label=(\alph*)]
		\item \label{prop-property-frechet-differentiable-a} Suppose $J_1, J_2: \calH \to \Real$ are Frech{\'e}t differentiable at $f \in \calH$ and $\alpha, \beta \in \Real$ are arbitrary. Then, $\alpha J_1 + \beta J_2$ is also Frech{\'e}t differentiable at $f \in \calH$, and 
		\begin{align*}
			\Diff \parens{\alpha J_1 + \beta J_2} \parens{f} = \alpha \Diff J_1 \parens{f} + \beta \Diff J_2 \parens{f}. 
		\end{align*}
		
		\item \label{prop-property-frechet-differentiable-b} (Chain rule) Suppose $J_1: \calH \to \Real$ is Frech{\'e}t differentiable at $f \in \calH$ and $J_2: \Real \to \Real$ is differentiable at $J_1 \parens{f}$. Then, $J_2 \circ J_1: \calH \to \Real$ is Frech{\'e}t differentiable at $f \in \calH$, and 
		\begin{align}
			\Diff \parens{J_2 \circ J_1} \parens{f} = J_2' \parens{J_1 \parens{f}} \Diff J_1 \parens{f}. 
		\end{align}
	\end{enumerate}
\end{proposition}

\begin{definition}[Fr{\'e}chet gradient]\label{def-frechet-gradient}
	Suppose $J: \calH \to \Real$ is Frech{\'e}t differentiable at $f \in \calH$. Since $\Diff J \parens{f}$ is a bounded linear map from $\calH$ to $\Real$, the Riesz-Fr{\'e}chet representation theorem \parencites[Fact 2.24 in][]{Bauschke2017-he} guarantees there exists a unique element $\nabla J \parens{f} \in \calH$ such that, for any $g \in \calH$, 
	\begin{align}\label{eq-frechet-grad}
		\Diff J \parens{f} \parens{g} = \innerp{g}{\nabla J \parens{f}}_{\calH}, 
	\end{align}
	and $\nabla J \parens{f}$ is called the \textit{Fr{\'e}chet gradient} of $J$ at $f$. If $J$ is Fr{\'e}chet differentiable on $\calH$, the \textit{Fr{\'e}chet gradient operator} is defined to be $\nabla J: \calH \to \calH, f \mapsto \nabla J \parens{f}$. 
\end{definition}

\begin{remark}
	Note that $\Diff J \parens{f}$ is a bounded linear map from $\calH$ to $\Real$, and belongs to the dual space of $\calH$, denoted by $\calH^*$. Since we have $\Diff J \parens{f} \parens{g} = \innerp{\nabla J \parens{f}}{g}_{\calH}$, the Riesz-Fr{\'e}chet representation theorem implies that $\norm{\Diff J \parens{f}}_{\calH^*} = \norm{\nabla J \parens{f}}_{\calH}$, where $\norm{}_{\calH^*}$ denotes the norm of the dual space $\calH^*$. 
\end{remark}

We now extend Definition \ref{def-frechet-differentiable} to higher orders. 

\begin{definition}[Higher-order Fr{\'e}chet differentiability and derivatives]\label{def-frechet-derivative-higher-order}
	Higher-order Fr{\'e}chet differentiability and derivatives are defined inductively. 
	
	In particular, the map $J$ is said to be \textit{twice Fr{\'e}chet differentiable at $f \in \calH$} if $J$ itself is Fr{\'e}chet differentiable at $f \in \calH$ and the map $\Diff J \parens{f}: \calH \to \Real$ is also Fr{\'e}chet differentiable at $f \in \calH$. The \textit{second Fr{\'e}chet derivative} of $J$ at $f \in \calH$, denoted by $\Diff^2 J \parens{f}$, is an operator from $\calH$ to $\calB \parens{\calH, \Real}$, that satisfies 
	\begin{align}\label{eq-2ndfrederivative}
		\lim_{\substack{\norm{g}_{\calH} \to 0 \\ g \neq 0}} \frac{\norm[\big]{\Diff J \parens{f + g} - \Diff J \parens{f} - \Diff^2 J \parens{f} \parens{g} }_{\calH^*}}{\norm{g}_{\calH}} = 0, 
	\end{align}
	where $\norm{\,\cdot\,}_{\calH^*}$ denotes the norm of the dual space of $\calH$. The map $J$ is said to be \textit{twice Fr{\'e}chet differentiable on $\calH$} if it is twice Fr{\'e}chet differentiable at all $f \in \calH$. 
	
	Suppose $J$ is twice Fr{\'e}chet differentiable on $\calH$. The \textit{second-order Fr{\'e}chet gradient operator}, denoted by $\nabla^2 J$, is a bounded linear operator that maps from $\calH$ to $\calB \parens{\calH, \calH}$ and satisfies 
	\begin{align*}
		\Diff^2 J \parens{f} \parens{g} \parens{h} = \innerp{h}{ \nabla^2 J \parens{f} \parens{g}}_{\calH}, \qquad \text{ for all } f, g, h \in \calH. 
	\end{align*}
	In other words, $\nabla^2 J \in \calB \parens{\calH, \calB \parens{\calH, \calH}}$ and $\nabla^2 J \parens{f} \in \calB \parens{\calH, \calH}$ for all $f \in \calH$. 
\end{definition}

%\begin{remark}
%	By Proposition 5.1.17 in \textcite{Denkowski2013-ke}, there exists an isometric isomorphism between $\calB \parens{\calH, \calB \parens{\calH, \calH}}$ and $\calB \parens{\calH \times \calH, \calH}$. Let $\Phi$ denote this isometric isomorphism. Then, $\Phi \parens{\nabla^2 J}$ is a map from $\calH \times \calH$ to $\calH$ such that $\Phi \parens{\nabla^2 J} \parens{f, \, g} = \nabla^2 J \parens{f} \parens{g}$, for all $f, g \in \calH$. 
%\end{remark}


\section{Bochner Integral}\label{section-bochner-integral}

In this section, we present the definition of the Bochner integral, which is the extension of the Lebesgue integral of real-valued functions to the integral of functions taking values in a Banach space. We also present some properties of the Bochner integral that we have used in the dissertation (in particular, in Chapter \ref{ch2} and \ref{ch3}). All materials of this section come from Appendix A.5.3 in \textcite{Steinwart2008-tn} and Section 3.10 \textcite{Denkowski2013-ke}. 

Throughout this section, let $\calE$ be a Banach space whose norm is denoted by $\norm{}_{\calE}$, and $\parens{\calX, \Sigma, \mu}$ be a $\sigma$-finite measure space (note that this $\mu$ differs from the one in the definition of finite-dimensional and kernel exponential families in Chapter \ref{ch2}). We first define the simple function (Definition \ref{def-simple-fun}) and the measurable function (Definition \ref{def-measurable-fun}) in the Banach space setting  and then define the Bochner $\mu$-integral (Definition \ref{def-bochner-integral}). 

\begin{definition}[$\calE$-valued simple function]\label{def-simple-fun}
	A function $s: \calX \to \calE$ is said to be an \textit{$\calE$-valued simple function} if there exist $e_1, \cdots, e_n \in \calE$ and $A_1, \cdots, A_n \in \Sigma$ such that 
	\begin{align*}
		s \parens{x} = \sum_{i=1}^n \indic_{A_i} \parens{x} e_i, \qquad \text{ for all } x \in \calX, 
	\end{align*}
	where $\indic_A$ is the indicator function of the set $A$, and is equal to 1 if $x \in A$ and to 0, otherwise. 
\end{definition}

\begin{definition}[$\calE$-valued measurable function]\label{def-measurable-fun}
	A function $f: \calX \to \calE$ is said to be an \textit{$\calE$-valued measurable function} if there exists a sequence of $\calE$-valued simple functions, $\sets{s_n}_{n \in \Natural}$, such that
	\begin{align}
		\lim_{n \to \infty} \norm{f \parens{x} - s_n \parens{x}}_{\calE} = 0
	\end{align}
	holds for all $x \in \calX$. 
\end{definition}

\begin{definition}[Bochner $\mu$-integral]\label{def-bochner-integral}
	An $\calE$-valued measurable function $f: \calX \to \calE$ is said to be \textit{Bochner $\mu$-integrable} if there exists a sequence of $\calE$-valued simple functions, $\sets{s_n}_{n \in \Natural}$, such that
	\begin{align}
		\lim_{n \to \infty} \int_{\calX} \norm{s_n \parens{x} - f \parens{x}}_{\calE} \ \diff \mu \parens{x} = 0. 
	\end{align}
	In this case, the limit
	\begin{align*}
		\int_{\calX} f \parens{x} \diff \mu \parens{x} := \lim_{n \to \infty} \int_{\calX} s_n \parens{x} \diff \mu \parens{x}
	\end{align*}
	exists and is called the \textit{Bochner $\mu$-integral} of $f$. 
\end{definition}

A criterion to check the Bochner $\mu$-integrability is the following. 

\begin{proposition}\label{prop-bochner-integrability}
	A measurable function $f: \calX \to \calE$ is Bochner $\mu$-integrable if and only if $\int_{\calX} \norm{f \parens{x}}_{\calE} \diff \mu \parens{x} < \infty$. 
\end{proposition}

Finally, we look at some properties of Bochner $\mu$-integral we use. 

\begin{proposition}\label{prop-properties-bochcher-integral}
	The Bochner $\mu$-integral defined above has the following properties: 
	\begin{enumerate}[label=(\alph*)]
		\item \label{prop-properties-bochcher-integral-a} The Bochner $\mu$-integral is linear. 
		\item \label{prop-properties-bochcher-integral-b} If $f: \calX \to \calE$ is Bochner $\mu$-integrable, we have 
		\begin{align*}
			\norm[\bigg]{\int_{\calX} f \parens{x} \diff \mu \parens{x}}_{\calE} \le \int_{\calX} \norm{f \parens{x}}_{\calE} \diff \mu \parens{x}. 
		\end{align*}
		\item \label{prop-properties-bochcher-integral-c} Suppose $\calE'$ is another Banach space. If $S: \calE \to \calE'$ is a bounded linear operator and $f: \calX \to \calE$ is Bochner $\mu$-integrable, then $S \circ f: \calX \to \calE'$ is also Bochner $\mu$-integrable. In this case, the integral commutes with $S$, that is, 
		\begin{align*}
			S \parens[\bigg]{\int_{\calX} f \parens{x} \diff \mu \parens{x}} = \int_{\calX} \parens{S \circ f} \parens{x} \, \diff \mu \parens{x}. 
		\end{align*}
	\end{enumerate}
\end{proposition}


\section{Partial Derivatives of a Kernel Function}\label{section-kernel-derivative}

In this section, we discuss the partial derivatives of the kernel function associated with a RKHS and their reproducing property. We follow the development in Section 4.3 in \textcite{Steinwart2008-tn} and the paper by \textcite{Zhou2008-jt}. Throughout this section, we let $\calX \subseteq \Real^d$ be an open set and $\Natural_0 := \Natural \cup \sets{0}$. 

We first consider a real-valued function $f: \calX \to \Real$. The function $f$ is said to be \textit{$m$-times continuously differentiable} if, for all $\alpha := \parens{\alpha_1, \cdots, \alpha_d}^\top \in \Natural_0^d$ with $\abs{\alpha} := \sum_{i=1}^d \alpha_i \le m$ and all $x \in \calX$, 
\begin{align*}
	\partial^\alpha f \parens{x} 
	%= \partial_1^{\alpha_1} \cdots \partial_d^{\alpha_d} f \parens{x} 
	= \frac{\partial^{\abs{\alpha}}}{\partial u_1^{\alpha_1} \cdots \partial u_d^{\alpha_d}} f \parens{u} \bigg\vert_{u=x}, 
\end{align*}
exists, where $u := \parens{u_1, \cdots, u_d}^\top \in \calX$. 

%\begin{lemma}[Differentiability of feature maps (Lemma 4.34 in \cite{Steinwart2008-tn})]\label{derivative.reproducing}
%	Let $\calX \subseteq \Real^d$ be an \textit{open} subset, $k$ be a kernel on $\calX$, $\calH$ be the RKHS of $k$, and $\Phi: \calX \to \calH$ be a feature map of $k$. Let $i \in \sets{1, \cdots, d}$ be an index such that the mixed partial derivative $\partial_i \partial_{i+d} k$ of $k$ with respect to the coordinates $i$ and $i+d$ exists and is continuous, where. Then the partial derivative $\partial_i \Phi$ of $\Phi: \calX \to \hil$ with respect to the $i$-th coordinate exists, is continuous, and for all $x, x \in \calX$, we have
%	\begin{align}
%		\innerp{\partial_i \Phi \parens{x}}{\partial_{i} \Phi \parens{x'}}_{\calH} = \partial_i \partial_{i+d} k \parens{x, x'} = \partial_{i+d} \partial_i k \parens{x, x'}. 
%	\end{align}
%\end{lemma}

We then define the $m$-times continuous differentiability of the kernel function $k$. 

\begin{definition}[$m$-times continuous differentiability of a kernel function]
	Let $k: \calX \times \calX \to \Real$ be a kernel function and $m \in \Natural$. We say $k$ is \textit{$m$-times continuously differentiable} if $\partial^{\alpha, \alpha}k: \calX \times \calX \to \Real$ exists and is continuous for all $\alpha := \parens{\alpha_1, \cdots, \alpha_d} \in \Natural_0^d$ with $\abs{\alpha} := \sum_{i=1}^d \alpha_i \le m$, where 
	\begin{align*}
		\partial^{\alpha, \alpha} k \parens{x, y} 
		%= \partial_1^{\alpha_1} \cdots \partial_d^{\alpha_d} \partial_{1+d}^{\alpha_1} \cdots \partial_{2d}^{\alpha_d} k \parens{x, y} 
		= \frac{\partial^{2\abs{\alpha}}}{\partial u_1^{\alpha_1} \cdots \partial u_d^{\alpha_d} \partial v_1^{\alpha_1} \cdots \partial v_d^{\alpha_d}} k \parens{u, v}\bigg\vert_{u=x, v=y}, \qquad \text{ for all } x, y \in \calX. 
	\end{align*}
\end{definition}

%\begin{remark}
%	\begin{enumerate}
%	\item From the lemma above, it is observed that $\partial_i \partial_{i+d} k$ is a kernel on $\calX \times \calX$ with the feature map $\partial_i \Phi$. 
%	\item Assuming that $\partial_i \partial_{i+d} \partial_j \partial_{j+d} k$ exists and is continuous, applying the lemma above iteratively one can show that $\partial_j \partial_i \Phi$ exists, is continuous, and is a feature map of the kernel $\partial_i \partial_{i+d} \partial_j \partial_{j+d} k$. 
%	\end{enumerate}
%\end{remark}

The partial derivative of $k$ is an element in $\calH$ and has the reproducing property as $k$ does, as the following proposition states. 

\begin{proposition}[Partial derivatives of kernels and its reproducing property]\label{prop-reproducing-derivative}
	Let $\calH$ be a RKHS with the kernel function $k: \calX \times \calX \to \Real$, and assume $k$ is $m$-times continuously differentiable on $\calX$. Then, 
	\begin{enumerate}[label=(\alph*)]
		\item we have 
		\begin{align}
			\partial^{\alpha} k \parens{x, \,\cdot\,} 
			%= \partial_1^{\alpha_1} \cdots \partial_d^{\alpha_d} k \parens{x, \,\cdot\,} 
			= \frac{\partial^{\abs{\alpha}}}{\partial u_1^{\alpha_1} \cdots \partial u_d^{\alpha_d}} k \parens{u, \,\cdot\,}\bigg\vert_{u=x} \in \calH
		\end{align}
		where $u := \parens{u_1, \cdots, u_d}^\top \in \calX$, and 
		\item every $f \in \hil$ is $m$-times continuously differentiable, and, for all $\alpha \in \Natural_0^d$ with $\abs{\alpha} \le m$ and all $x \in \calX$, the partial derivative reproducing property holds, i.e., 
		\begin{align}
			\partial^{\alpha} f \parens{x} = \innerp{\partial^{\alpha} k \parens{x, \,\cdot\,}}{f}_{\calH}, \qquad \text{ for all } x \in \calX. %, \hspace{20pt} \text{and}\hspace{20pt}  \abs{\partial^{\alpha} f(x)} \le \norm{f}_{\hil} \cdot \sqrt{\partial^{\alpha, \alpha} k\pare{x, x}}. 
		\end{align}
		In particular, we have $\partial^{\alpha, \alpha} k \parens{x, y} = \innerp{\partial^{\alpha} k \parens{x, \,\cdot\,}}{\partial^{\alpha} k \parens{y,\,\cdot\,}}_{\calH}$ for all $x, y \in \calX$. 
	\end{enumerate}
\end{proposition}


\section{Some Theories on Bounded Linear Operators}\label{section-bdd-linear-operator}

In our development in Chapters \ref{ch2} and \ref{ch3}, we have used some theories on the bounded linear operators in a Hilbert space. In this section, we aim to give a very brief overview of these theories. More information can be found in, for example, Chapters 13 and 16 in \textcite{Royden2018-cd} and Chapter VI in \textcite{Reed2012-xs}. 

Throughout this section, we let $\calH$ be a real Hilbert space with the inner product $\innerp{\,\cdot\,}{\,\cdot\,}_{\calH}$ and the norm $\norm{}_{\calH}$, and assume $\calH$ is separable, meaning that it contains a dense countable subset. Due to the separability of $\calH$, it admits a countable orthonormal basis. 

\begin{definition}[Linear operator]
	An operator $C: \calH \to \calH$ is said to be \textit{linear} if, for any $f, g \in \calH$ and any $\alpha, \beta \in \Real$, we have 
	\begin{align*}
		C \parens{\alpha f + \beta g} = \alpha C f + \beta C g. 
	\end{align*}
\end{definition}

\begin{definition}[Bounded operator]
	A linear operator $C: \calH \to \calH$ is said to be \textit{bounded} if there exists a constant $M \ge 0$ such that 
	\begin{align*}
		\norm{C f}_{\calH} \le M \cdot \norm{f}_{\calH}, \qquad \text{ for all } f \in \calH. 
	\end{align*}
	The infimum of all such $M$ is called the \textit{operator norm} of $C$ and is denoted by $\norm{C}$. 
\end{definition}

\begin{definition}[Adjoint and self-adjoint operators]
	Let $C: \calH \to \calH$ be a bounded linear operator. Then, the \textit{adjoint} of $C$ is the bounded linear operator $C^*: \calH \to \calH$ satisfying 
	\begin{align*}
		\innerp{Cf}{g}_{\calH} = \innerp{f}{C^*g}_{\calH}, \qquad \text{ for all } f, g \in \calH. 
	\end{align*}
	The adjoint $C^*$ exists and is unique. 
	
	The operator $C$ is said to be \textit{self-adjoint}, if $C = C^*$, that is, $\innerp{Cf}{g}_{\calH} = \innerp{f}{Cg}_{\calH}$ for all $f, g \in \calH$. 
\end{definition}

\begin{definition}[Positive semidefinite and definite operators]
	Let $C: \calH \to \calH$ be a self-adjoint bounded linear operator. Then, $C$ is said to be \textit{positive semidefinite} if, for all $f \in \calH$, $\innerp{f}{Cf}_{\calH} \ge 0$, and to be \textit{positive definite} if, for all $f \in \calH \backslash \sets{0}$, $\innerp{f}{Cf}_{\calH} > 0$. 
\end{definition}

\begin{definition}[Compact operator]
	A bounded linear operator $C: \calH \to \calH$ is said to be \textit{compact}, if, for every bounded sequence $\sets{f_n}_{n \in \Natural}$ in $\calH$, $\sets{C f_n}_{n \in \Natural}$ has a subsequence that converges in $\calH$ (with respect to the norm $\norm{}_{\calH}$). 
\end{definition}

Other equivalent ways of defining a compact operator can be found, for example, in Section 16.5 in \textcite{Royden2018-cd} and Section VI.5 in \textcite{Reed2012-xs}. 

\begin{definition}[Finite rank operator]
	A bounded linear operator $C: \calH \to \calH$ is said to be of \textit{finite rank} if its range is finite-dimensional. That is, every element in $\range \parens{C}$ can be written as $Cf = \sum_{i=1}^m \alpha_i g_i$, for some $f \in \calH$, some fixed family $\sets{g_i}_{i=1}^m \subset \calH$, some $\alpha_1, \cdots, \alpha_m \in \Real$, and $m < \infty$. 
\end{definition}

The relationship between the compact operator and the finite rank operator is given in the following proposition. 

\begin{proposition}
	Every finite rank operator is compact. In addition, every compact operator on $\calH$ is the norm limit of a sequence of finite rank operators. 
\end{proposition}

In addition, we have the following characterization when the identity operator is compact. 

\begin{proposition}
	The identity operator $I: \calH \to \calH$ is compact if and only if $\calH$ is finite-dimensional. 
\end{proposition}

We also have the following result characterizing the relationship between the invertibility and the compactness of a bounded linear operator. 

\begin{proposition}\label{prop-invertible-compact-infdim}
	Let $\calH$ be infinite-dimensional and $C: \calH \to \calH$ be a bounded linear operator. If $C$ is compact, it is \emph{not} invertible. 
\end{proposition}

In the development in Chapter \ref{ch3}, we have repeatedly used the following Hilbert-Schmidt theorem, which is an extension of the eigen-decomposition of a real symmetric matrix. 

\begin{theorem}[Hilbert-Schmidt theorem]\label{theorem-hilbert-schmidt}
	Let $C: \calH \to \calH$ be a self-adjoint compact operator. There exists an orthonormal basis $\sets{\psi_{\nu}}_{\nu=1}^R$ for $\overline{\range \parens{C}}$, together with a monotonically non-increasing sequence of nonzero real numbers $\sets{\xi_{\nu}}_{\nu=1}^R$, such that $C \psi_{\nu} = \xi_{\nu} \psi_{\nu}$ for all $\nu = 1, \cdots, R$. In addition, the following identity holds 
	\begin{align*}
		C f = \sum_{\nu=1}^R \xi_{\nu} \innerp{f}{\psi_{\nu}}_{\calH} \psi_{\nu}, \qquad \text{ for all } f \in \calH. 
	\end{align*}
	If $C$ is of finite rank, then $R < \infty$ and $R$ is the rank of $C$; if $C$ is not of finite rank, $R = \infty$ and $\lim_{\nu \to \infty} \xi_{\nu} = 0$. 
\end{theorem}

In Theorem \ref{theorem-hilbert-schmidt}, the functions $\psi_{1}, \cdots, \psi_{R}$ are called the \textit{eigenfunctions} of $C$, and $\xi_{\nu}$ is called the \textit{eigenvalue} of $C$ associated with the eigenfunction $\psi_{\nu}$, for all $\nu = 1, \cdots, R$. 

Using Theorem \ref{theorem-hilbert-schmidt}, it is easy to see that, if $C$ is positive semidefinite, all its eigenvalues are positive. 

We next introduce two special classes of bounded linear operators, the trace class and the Hilbert-Schmidt class, that are of great importance in proving various properties of penalized and early stopping SM density estimators in Chapter \ref{ch2} and \ref{ch3}. 

We first introduce the trace of a bounded linear operator, based on which we define the trace class. 

\begin{definition}[Trace]
	Let $\sets{\psi_{\nu}}_{\nu=1}^{\infty}$ be an orthonormal basis of $\calH$. For any positive definite bounded linear operator $C: \calH \to \calH$, the \textit{trace} of $C$ is defined to be 
	\begin{align*}
		\trace \parens{C} := \sum_{\nu=1}^{\infty} \innerp{\psi_{\nu}}{C \psi_{\nu}}_{\calH}, 
	\end{align*}
	where the sum is independent of the choice of the orthonormal basis. 
\end{definition}

\begin{definition}[Trace class]
	A bounded linear operator $C: \calH \to \calH$ is said to be of \textit{trace class} if $\tr \parens{\abs{C}} < \infty$. 
\end{definition}

Then, we have the following properties of operators that are of trace class. 

\begin{proposition}
	Let $C: \calH \to \calH$ be of trace class. Then, $C$ is compact, and its operator norm, $\norm{C}$, and its trace, $\tr \parens{C}$, are related by $\norm{C} \le \tr \parens{C}$. 
\end{proposition}

We now turn to the Hilbert-Schmidt operator. 

\begin{definition}[Hilbert-Schmidt operator]
	Let $\sets{\psi_{\nu}}_{\nu=1}^{\infty}$ be an orthonormal basis of $\calH$. A bounded linear operator $C: \hil \to \hil$ is said to be \textit{Hilbert-Schmidt (HS)} if 
	\begin{align*}
		\sum_{\nu=1}^{\infty} \norm{C \psi_{\nu}}^2_{\calH} < \infty. 
	\end{align*}
\end{definition}

Still let $\sets{\psi_{\nu}}_{\nu=1}^{\infty}$ be an orthonormal basis of $\calH$. Given two HS operators $C_1, C_2: \hil \to \hil$, define the \textit{HS inner product} between them to be 
\begin{align*}
	\innerp{C_1}{C_2}_{\mathrm{HS}} = \sum_{\nu=1}^{\infty} \innerp{C_1 \psi_{\nu}}{C_2 \psi_{\nu}}_{\hil}, 
\end{align*}
and the \textit{HS norm} to be 
\begin{align*}
	\norm{C_1}_{\mathrm{HS}} := \sqrt{\innerp{C_1}{C_1}_{\mathrm{HS}}} = \sqrt{\tr \parens{C_1^* C_1}} = \parens[\Bigg]{\sum_{\nu=1}^{\infty} \norm{C_1 \psi_{\nu}}^2_{\hil}}^{1/2}. 
\end{align*}

We then have the following properties of HS operators. 

\begin{proposition}[Properties of HS operators]\label{prop-hs-property}
	The following properties of the HS operators hold: 
	\begin{enumerate}[label=(\alph*)]
		\item \label{prop-hs-property-a} Let $C: \hil \to \hil$ be a HS operator. Then, its operator norm $\norm{C}$, its trace $\trace \parens{C}$, and its HS norm $\norm{C}_{\mathrm{HS}}$ are related by $\norm{C} \le \norm{C}_{\mathrm{HS}} \le \trace \parens{C}$. 
		\item \label{prop-hs-property-b} The class of all HS operators with the inner product $\innerp{\,\cdot\,}{\,\cdot\,}_{\mathrm{HS}}$ forms a Hilbert space. 
		\item \label{prop-hs-property-c} Every HS operator is compact. 
	\end{enumerate}
\end{proposition}

